\subsection{Intervjuer}
För beskrivning av datainsamlingsmetoden, se avsnitt~\ref{subsubsec:datainsamling}.
\subsubsection{Lokförare och DAS-verktyg}
Den första intervjufrågan lyder \emph{\lq{}Vilka erfarenheter har du av DAS-applikationer idag?\rq{}}, och är direkt till forskningsfråga 1 som lyder \emph{\lq{}Hur ser benägenheten ut för lokförare att följa hastighetsrekommendationer från en DAS-applikation idag?\rq{}}.
Frågan syftade till att försöka bilda en uppfattning om hur stor kännedom lokförare idag har av digitala förarverktyg som stöttar denne i tågets framförande.
Samtliga respondenter hade erfarenhet av S-DAS-applikationer, och fyra av dem använde en DAS-applikation i deras dagliga arbete.
Verktyget, vilket det än må vara uppgavs vara väldigt uppskattat.
Särskilt uppskattade funktioner uppgavs i synnerhet vara (nära) realtidsuppdaterad operativ omlägesbild med framför- och bakomvarande tåg samt möjlighet till körorderuttag\footnote{En körorder består av lokförarens tågorder och körplan och är ett villkor för att genomföra tågfärd~\cite{TTJ2015}. Körorderuttag är hämtandet av körordet och sker i form av nedladdning direkt från Trafikverkets API eller som PDF från körordersystemet.}.
Att kunna planera sin körning krävdes enligt en respondent linjekännedom\footnote{Kännedom om en banas egenskaper såsom driftplatsers geografiska lägen, signalplaceringar, lutningsförhållanden eller andra särskilda egenskaper för en viss bana eller driftplats.}~för att kunna planera sin tågkörning och beslut om hastigheter eller inbromsningar fattas av förare själva.
Några uppgav att eco driving-funktionalitet som rekommenderade en viss hastighet fungerade dåligt och därmed inte följdes.
Information om banans STH, tillfälliga hastighetsnedsättningar, hastighetsnedsättningar kopplat till fordon har inte varit tillgängligt för dessa system och de har därför kunnat uppge hastigheter som är för höga och inte går att följa.
Andra funktioner kunde vara kontaktuppgifter, dokument, fordonsomlopp och personligt schema.
Respondenten som var godstågsförare uppgav att denne kört med DAS-verktyg på LIA\footnote{Lärande i arbete, vanligt förekommande på yrkeshögskolor.}, men att det var ovanligt på hans företag att lokförare körde med förarverktyg överhuvudtaget.
Tågorder och körplan skrevs ut vid orderläsning\footnote{Arbetsmoment där lokföraren läser in sig om särskilda trafiksäkerhetsförhållanden som gäller för dagen och om uppdateringar av regler förekommit sedan senaste arbetspasset}.
Några hade deltagit i pilotstudien som Trafikverket genomförde på Södra Stambanan och därmed fått testa C-DAS.
Dessa uppgav att det sällan fungerade, och att de hade svårt att förstå syftet med systemet.
\paragraph{}
\hspace{-11}
Överlag var intervjuobjekten positivt inställda till DAS-verktyg, någon uttryckte att det var \lq{}framtiden\rq{}~, att järnvägen självklart skulle gå i bräschen för hållbar transport och hoppades på ett snart införande.
Någon uttryckte en oro över att yrkesrollen skulle förändras till att vara en \lq{}DAS-vakt\rq{}~och ta bort \lq{}hantverket\rq{}~om DAS-verktyg fick för stor inverkan på arbetsutövandet istället för att lokföraren kör tåg på sättet denne själv anser bäst.
\subsubsection{Kommunikations- och informationsbehov}
Den andra intervjufrågan lyder \emph{\lq{}\ldots~vilken information är viktig/krävs för dig som förare för att du ska följa en hastighetsanvisning?\rq{}}.
Frågan är kopplad till forskningsfråga 2 som lyder \emph{\lq{}Hur väl tillgodoser befintlig standard för datautbyte inom C-DAS (SFERA) lokförares informationsbehov avseende efterlevnad av hastighetsanvisningar?\rq{}}.
\paragraph{}
\hspace{-11}Den tredje intervjufrågan lyder \emph{\lq{}Vilka förhållanden skulle kunna begränsa din förmåga att efterleva en hastighetsanvisning?\rq{}}, och är precis som intervjufråga 2 även kopplad till forskningsfråga 2.
Vid några tillfällen presenterades på nytt, parafraserat, i stil med \emph{Vad skulle kunna föranleda att du ringer till fjärren\footnote{\lq{}Fjärren\rq{} är benämning på Trafikverkets driftledningscentral(er) och kommer från ordet \lq{}fjärrtågklarerare\rq{}.}~för att meddela att du kör med sådan förutsättningar att du inte kan köra enligt din körplan?}~för att komma till kärnan av frågan som var att identifiera kommunikation som är relevant för C-DAS\@.


\subsubsection{SFERA utökat}

\paragraph{}
\hspace{-11}Datainsamling för den sista forskningsfrågan har skett genom dokumentanalys och presenteras